% Options for packages loaded elsewhere
\PassOptionsToPackage{unicode}{hyperref}
\PassOptionsToPackage{hyphens}{url}
\PassOptionsToPackage{dvipsnames,svgnames,x11names}{xcolor}
%
\documentclass[
  10pt,
]{article}
\usepackage{amsmath,amssymb}
\usepackage{iftex}
\ifPDFTeX
  \usepackage[T1]{fontenc}
  \usepackage[utf8]{inputenc}
  \usepackage{textcomp} % provide euro and other symbols
\else % if luatex or xetex
  \usepackage{unicode-math} % this also loads fontspec
  \defaultfontfeatures{Scale=MatchLowercase}
  \defaultfontfeatures[\rmfamily]{Ligatures=TeX,Scale=1}
\fi
\usepackage{lmodern}
\ifPDFTeX\else
  % xetex/luatex font selection
  % Use the default Latin Modern fonts supplied with TeX.

\fi
% Use upquote if available, for straight quotes in verbatim environments
\IfFileExists{upquote.sty}{\usepackage{upquote}}{}
\IfFileExists{microtype.sty}{% use microtype if available
  \usepackage[]{microtype}
  \UseMicrotypeSet[protrusion]{basicmath} % disable protrusion for tt fonts
}{}
\makeatletter
\@ifundefined{KOMAClassName}{% if non-KOMA class
  \IfFileExists{parskip.sty}{%
    \usepackage{parskip}
  }{% else
    \setlength{\parindent}{0pt}
    \setlength{\parskip}{6pt plus 2pt minus 1pt}}
}{% if KOMA class
  \KOMAoptions{parskip=half}}
\makeatother
\usepackage{xcolor}
\usepackage[margin=24mm]{geometry}
\setlength{\emergencystretch}{3em} % prevent overfull lines
\providecommand{\tightlist}{%
  \setlength{\itemsep}{0pt}\setlength{\parskip}{0pt}}
\setcounter{secnumdepth}{-\maxdimen} % remove section numbering
\usepackage{microtype}
\setlength{\emergencystretch}{3em}
\widowpenalty=10000
\clubpenalty=10000
\ifLuaTeX
  \usepackage{selnolig}  % disable illegal ligatures
\fi
\usepackage{bookmark}
\IfFileExists{xurl.sty}{\usepackage{xurl}}{} % add URL line breaks if available
\urlstyle{same}
\hypersetup{
  pdftitle={Ternary Measurements for Paired Local Recovery},
  colorlinks=true,
  linkcolor={blue},
  filecolor={Maroon},
  citecolor={Blue},
  urlcolor={blue},
  pdfcreator={LaTeX via pandoc}}

\title{Ternary Measurements for Paired Local Recovery}
\author{Chuhang Ding (Peking University)}
\date{25 November 2026}

\begin{document}
\maketitle

\subsection{Abstract}\label{abstract}

An input to paired local recovery has one scalar for each ordered pair
of length-\(q\) words over an \(r\)-letter alphabet. At the \(t\)-th
evaluation point, an output can read the inputs whose first word has the
prescribed letter at position \(t\). All outputs share a single linear
measurement matrix. We determine its minimum rank at two degree bounds
for the binary indexing alphabet. At degree zero over the rationals the
answer is \((4^q-1)/3\). A three-color tensor code with ternary rows
attains the converse of Woolsey, Chen and Ji for the centralized version
of their hypercuboid assignment. An adjacent interchange proves the
optimality of their proposed receiver order for unequal alphabet sizes
and gives the generalized independence number. For every fixed even
\(r\), a construction from perfect matchings attains their limiting
measurement fraction \((r-1)/(2r-1)\). At degree \(q-2\), the binary
minimum is \(q+\lceil q/2\rceil+2\) over every field with \(q\) distinct
nonzero evaluation points. A binary-valued measurement matrix attains
the minimum. The lower bound uses a degree-one interpolation dual and a
count of exceptional incidences on cube edges. At both binary endpoints,
imposing a direct sum over the coordinate-swap orbits strictly increases
the rank for \(q\ge2\), even when the measurements inside each orbit are
unrestricted.

\subsection{1. The model and the two binary
endpoints}\label{the-model-and-the-two-binary-endpoints}

Let \([r]=\{0,\ldots,r-1\}\) and let \(\Omega=[r]^q\times[r]^q\), with
\(D=r^{2q}\). Take a field \(F\) and distinct nonzero evaluation points
\(\lambda_1,\ldots,\lambda_q\in F\). Write \(v=(a_v,b_v)\) for a pair of
words. For \(u\in\Omega\), the side information at slot \(t\) consists
of the input coordinates in \[
G_{u,t}=\{v\in\Omega:a_v(t)=b_u(t)\}.
\] The sender knows the whole input \(z\in F^\Omega\). It broadcasts
\(Cz\) for a fixed matrix \(C\) with \(D\) columns. A feasible
degree-\(d\) scheme has a matrix \(P\in F[X]^{D\times D}\) of entrywise
degree at most \(d\) and \(P(0)=I_D\). For every \(u,t\) the value
\([P(\lambda_t)z]_u\) must be recoverable from \(Cz\) and
\(z|_{G_{u,t}}\) for every \(z\in F^\Omega\). Neither \(C\) nor \(P\) is
required to respect any permutation of \(\Omega\). Over \(\mathbb Q\) a
ternary matrix has entries in \(\{-1,0,1\}\). An independent subset of
its rows preserves everything the broadcast reveals, so its rank is the
minimum number of rows needed to transmit that row space.

We write \(m_F(q,r,d)\) for the minimum rank with arbitrary entries in
\(C\) and \(m_{\rm ter}(q,r,d)\) for the minimum using ternary rows over
\(\mathbb Q\). The labels are fixed in these quantities. The formulas
below hold for every allowed choice. A linear functional is recoverable
precisely when it vanishes on the common kernel of the broadcast and the
side information. Consequently each row of \(P(\lambda_t)\) lies in the
sum of \(\operatorname{Row}(C)\) and the span of the coordinate
functionals indexed by \(G_{u,t}\). This equivalence also allows
arbitrary decoding functions.

The paired label array, coordinate swaps and polynomial replacement rows
arise in the common auxiliary construction of {[}1, §4.1{]}. At degree
zero we have \(P=I_D\), and the receiver \((t,b)\) knows all inputs with
first letter \(b\) at \(t\) and demands all inputs with second letter
\(b\) there. These are the file and function assignments of Woolsey,
Chen and Ji {[}2, §III-B{]}, with a single scalar per file-function
pair. Their communication protocol further requires each sender to form
its transmissions from locally available data. Our sender can combine
inputs anywhere in the array. Receiver-order converses for such
groupcast problems also appear in {[}3, §5{]}. The moment obstruction
for a fixed common space in the polynomial model appears in {[}4,
§5.1{]}. We use its elementary interpolation form at degree \(q-2\).

\textbf{Theorem 1.} Over \(\mathbb Q\) the binary indexing alphabet
satisfies \[
\begin{aligned}
 m_{\mathbb Q}(q,2,0)=m_{\rm ter}(q,2,0)
   &=\frac{4^q-1}{3} &&(q\ge1),\\
 m_{\mathbb Q}(q,2,q-2)=m_{\rm ter}(q,2,q-2)
   &=q+\left\lceil\frac q2\right\rceil+2 &&(q\ge2).
\end{aligned}
\] The second formula holds for \(m_F\) over every field containing the
required labels. Its attaining measurement matrix has entries \(0,1\).
The requirement on the labels rules out the field \(\mathbf F_2\) itself
when \(q\ge2\). Fields of characteristic two with sufficiently many
elements are included.

\textbf{Theorem 2.} If \(r\ge2\) is fixed and even, then \[
 \lim_{q\to\infty}\frac{m_{\mathbb Q}(q,r,0)}{r^{2q}}
 =\lim_{q\to\infty}\frac{m_{\rm ter}(q,r,0)}{r^{2q}}
 =\frac{r-1}{2r-1}.
\] For all \(q\ge1\) there are explicit ternary codes of rank at most \[
 \begin{cases}
 \dfrac{r-1}{2r-1}\left(r^{2q}-(r/2)^q\right),
       &r\text{ even},\\[3pt]
 \dfrac{r}{2r+1}\left(r^{2q}-((r+1)/2)^q\right),
       &r\text{ odd}.
 \end{cases}
\]

The two assertions of Theorem 1 coincide at \(q=2\), where five
measurements suffice. At \(q=3\) their ranks are \(21\) and \(7\). The
saving against orbit-by-orbit measurement has a different cause at the
two degrees. At degree zero, sums indexed by a color in a tensor basis
can serve all slots simultaneously. At degree \(q-2\), the
two-dimensional space of affine polynomials permits two different slots
to share one measurement direction.

\subsection{2. A published converse for degree
zero}\label{a-published-converse-for-degree-zero}

We give the receiver-order count of Woolsey, Chen and Ji {[}2, Theorem
2, equation (12) and Appendix C{]} to show its applicability to a
central sender and to specify its finite value. Order the \(rq\)
receivers \((t,c)\) first by increasing letter \(c\), then by increasing
slot \(t\). Select the message \((a,b)\) if the earliest receiver
demanding it precedes the earliest receiver holding it. If a kernel
vector of \(C\) is supported on the selected messages, visit the
receivers in increasing order. At the first demand for a selected
message, every selected message in this receiver's side information has
an earlier first demand. By induction all its side information vanishes.
The zero broadcast then forces all of its demanded coordinates to
vanish. Thus all selected columns of \(C\) are independent. The
analogous recovery with the complementary inputs revealed is the usual
information-theoretic permutation argument.

Fix the earliest receiver \((t,c)\) demanding a selected message and put
\(j=r-c\). Among the \(b\)-words, the first \(t-1\) positions must use
larger letters than \(c\), position \(t\) is fixed to \(c\), and the
later positions may also use \(c\). There are \((j-1)^{t-1}j^{q-t}\)
choices. The \(a\)-word must avoid holding receivers through \((t,c)\)
itself, giving \((j-1)^t j^{q-t}\) choices. The number of selected
messages is therefore \[
 B(q,r)=
 \sum_{j=2}^r\sum_{t=1}^q(j-1)^{2t-1}j^{2(q-t)}
 =
 \sum_{j=2}^r\frac{j-1}{2j-1}
       \left(j^{2q}-(j-1)^{2q}\right).
 \tag{1}
\] Using the difference between successive coefficients
\((j-1)/(2j-1)\), equation (1) yields \[
 \frac{m_F(q,r,0)}{r^{2q}}\ge
 \frac{r-1}{2r-1}
   -\sum_{j=1}^{r-1}\frac{(j/r)^{2q}}{4j^2-1}.
 \tag{2}
\] For \(r=2\) this gives the unnormalized lower bound \((4^q-1)/3\).
For fixed \(r\) the limiting lower bound is \((r-1)/(2r-1)\).

\textbf{Proposition 2.1 (optimal receiver order).} Allow slot \(i\) to
have \(r_i\) letters. Among all permutations of the receivers, a rule
that always chooses a slot with the most unvisited letters maximizes the
number of selected messages. With \(S(j)=\prod_i\min(r_i,j)^2\), the
maximum is \[
       \sum_{j=2}^{\max_i r_i}\frac{j-1}{2j-1}\bigl(S(j)-S(j-1)\bigr).
\] This proves the optimality suggested by simulations in {[}2, Remark
5{]} for the ordering rule used for \(L_{P2}\) in their Appendix
D.\footnote{In arXiv:2008.00581v1, equations (15) and (30) print the
  outside factor as \(\prod_{j<i}x_j^2\) for nonincreasing side lengths
  \(x_1,\ldots,x_q\). When the first \(i\) slots decrease together, the
  untouched slots give \(\prod_{j>i}x_j^2\). For \((x_1,x_2)=(2,2)\) the
  printed equation (15) gives \(5/4\), while its stated ordering rule
  and the homogeneous equation (12) give the normalized score \(5/16\).}
By {[}3, Theorem 5, equation (43){]}, the maximum is also the
generalized independence number of this unit-message groupcast instance,
after removing demands already available to their requesting receivers.

\emph{Proof.} Suppose \(s_i\) letters of slot \(i\) remain. Selecting a
receiver in slot \(t\) contributes \((s_t-1)\prod_{i\ne t}s_i^2\).
Indeed its demand letter is fixed, its holder letter must be different,
and all other letter pairs are free among the unvisited letters. For
consecutive choices in distinct slots \(t,u\), put \(x=s_t\) and
\(y=s_u\). Later contributions agree in either order. The score for
order \(t,u\) minus the score for \(u,t\) is
\(\prod_{i\notin\{t,u\}}s_i^2\) times \[
 (x-1)y^2+(y-1)(x-1)^2-(y-1)x^2-(x-1)(y-1)^2=x-y .
\] Tag each step by its slot's count just before that choice. Within a
slot the tags strictly decrease. Any adjacent inversion therefore uses
distinct slots. Swapping it leaves the tags unchanged and does not
decrease the score. Sorting all tags into nonincreasing order gives the
claimed rule. At tag \(j\) the slots with \(r_i\ge j\) decrease from
\(j\) to \(j-1\) in turn. Summing their contributions as a geometric
series, as in (1), gives \((j-1)(S(j)-S(j-1))/(2j-1)\). Summing over
\(j\) proves the formula. \(\square\)

\subsection{3. Centralized tensor
encoders}\label{centralized-tensor-encoders}

Write \(V\) for the one-slot dual space with coordinate basis
\(e_{ab}\), \(a,b\in[r]\), and put
\(K_b=\operatorname{span}\{e_{bc}:c\in[r]\}\). A subspace
\(U\subseteq V\) suffices in one slot if \(e_{ab}\in U+K_b\) for every
\(a,b\).

\textbf{Lemma 3 (selection by total color).} Partition a basis of \(V\)
into classes \((\mathcal B_g)_{g\in G}\) indexed by a finite abelian
group. Suppose \(A\subseteq G\) has the property that the classes with
colors in \(A+h\) span a sufficient space for every \(h\in G\).
Transmitting all tensor basis rows in \(V^{\otimes q}\) whose total
color is in a fixed translate \(A+h_0\) gives a degree-zero code. Its
rank is the sum over \(A+h_0\) of the \(q\)-fold convolution of
\(g\mapsto |\mathcal B_g|\). The tensor rows are ternary when the
one-slot basis is ternary.

\emph{Proof.} Fix receiver \((t,b)\) and a basis tensor on the other
slots of total color \(h\). At the remaining slot write \(e_{ab}\) as a
sum of a vector in the span of the colors \(A+h_0-h\) and a vector in
\(K_b\). After inserting these vectors at slot \(t\), the first part is
a linear combination of broadcast rows and the second uses only entries
known to the receiver. Expanding standard coordinate vectors on the
other slots in the tensor basis proves recovery. Every transmitted
vector belongs to one basis, so the row count gives its rank. Products
of entries preserve ternarity. \(\square\)

For the binary case take \(G=\mathbb Z/3\mathbb Z\) and the following
four vectors, grouped by color: \[
 \mathcal B_0=\{e_{00}+e_{10},\ e_{01}+e_{11}\},\qquad
 \mathcal B_1=\{e_{01}+e_{10}\},\qquad
 \mathcal B_2=\{e_{01}-e_{10}\}.
\] The last two vectors span the off-diagonal coordinates, and the first
two vectors project to independent diagonal coordinates. With side
information \(K_1\), each of \(e_{01}+e_{11}\), \(e_{01}+e_{10}\) and
\(e_{01}-e_{10}\) gives \(e_{01}\). With \(K_0\), use \(e_{00}+e_{10}\),
\(e_{01}+e_{10}\) and \(-(e_{01}-e_{10})\) to obtain \(e_{10}\). Each
self-matching coordinate is in its own side information. Thus every
color alone suffices. Transmit the tensor rows of total color one. The
single-slot multiplicities satisfy \((2,1,1)=\mathbf1+\delta_0\) on
\(G\). Convolving \(q\) times gives \((4^q-1)/3\) at each nonzero color
for every \(q\ge1\). Indeed,
\((\mathbf1+\delta_0)^{*q}=((4^q-1)/3)\mathbf1+\delta_0\). Together with
(1) this proves the degree-zero binary formula.

This binary indexing construction also gives the unrestricted minimum
over any field of size greater than two. Use \(e_{01}+\theta e_{10}\) in
place of the last vector for \(\theta\ne0,1\). Each color still suffices
and the four local rows are independent. Over the rationals the original
choice \(\theta=-1\) supplies ternary rows. Without polynomial labels
this central communication code works for arbitrary \(q\) over such a
field. If the polynomial model is retained, its assumption on label
availability still applies.

\subsubsection{3.1 Matchings for every
alphabet}\label{matchings-for-every-alphabet}

Put \(k=r-1\) for even \(r\) and \(k=r\) for odd \(r\), and set
\(p=2k+1\). Partition the edges of the complete graph on \(r\) letters
into \(k\) matchings. For even \(r\), use vertices
\(\{\infty\}\cup\mathbb Z/(r-1)\mathbb Z\). The matching with index
\(j\) consists of \(\{\infty,j\}\) and \(\{j+i,j-i\}\) for
\(1\le i\le(r-2)/2\). Each finite edge has a unique midpoint because
\(r-1\) is odd. For odd \(r\), take the perfect matchings for \(r+1\)
vertices and remove the extra vertex. Denote their common size by
\(m=\lfloor r/2\rfloor\).

There are \(k\) cyclic intervals of length \(k\) containing zero in
\(G=\mathbb Z/p\mathbb Z\). Assign a different interval to each
matching. For an edge \(\{a,b\}\) of the matching assigned
\(I=[j,j+k-1]\), give its rows \(e_{ab}+e_{ba}\) and \(e_{ab}-e_{ba}\)
colors \(j-1\) and \(j+k\), respectively, ordering \(a<b\). Give color
zero to each column sum \(H_b=\sum_a e_{ab}\). The edge rows span all
off-diagonal coordinates, and the column sums have independent diagonal
parts, so these rows form a ternary basis.

We check that \emph{every} cyclic interval of length \(k\) spans a
sufficient space. The complement of the two colors \(j-1,j+k\) consists
of paths on \(k\) and \(k-1\) residues. Accordingly \(I\) is the unique
interval missing both colors of the assigned edge. An interval omitting
zero sees at least one row of every edge. At receiver \(b\), either row
gives \(e_{ab}\) since \(e_{ba}\in K_b\). An interval containing zero
can miss only the edges in its assigned matching. At most one such edge
meets \(b\). Subtracting \(e_{bb}\) and all the already recovered
coordinates in column \(b\) from \(H_b\) recovers the missing
coordinate.

Take \(j=-k+1,\ldots,0\) for the interval starts. The left endpoints
outside the intervals exhaust \(-k,\ldots,-1\) and the right endpoints
exhaust \(1,\ldots,k\). Zero therefore has \(r\) basis rows and each
nonzero color has \(m\). The color multiplicity vector is
\(m\mathbf1+(r-m)\delta_0\), of total mass \(r^2\). Its \(q\)-fold
convolution has value \([r^{2q}-(r-m)^q]/p\) at each nonzero color.
Applying Lemma 3 with an interval of length \(k\) omitting zero gives
the finite ranks in Theorem 2. When \(r\) is even, dividing by
\(r^{2q}\) and using (2) proves the claimed exact limit. For odd
\(r\ge3\) this construction leaves a separation of \(1/(4r^2-1)\)
between the lower and upper limiting fractions.

\subsection{\texorpdfstring{4. Affine duality at degree
\(q-2\)}{4. Affine duality at degree q-2}}\label{affine-duality-at-degree-q-2}

For the rest of this section the indexing alphabet is binary, \(q\ge2\),
and \(F\) has the required labels. Let \(E=F[X]_{\le1}\). We first
justify the fixed-matrix criterion used on both sides of the rank bound.

\textbf{Lemma 4.} A matrix \(C\) admits a reconstruction matrix of
degree at most \(q-2\) if and only if the following holds for every
output \(u\). Each \(h\in F^\Omega\otimes E\) with \(Ch(X)=0\) and
\(h_v(\lambda_t)=0\) for \(v\in G_{u,t}\) satisfies \(h_u(0)=0\).

\emph{Proof.} Write \(K_t\) for the input subspace killed by \(C\) and
zero on \(G_{u,t}\). The proposed row
\(e_u^\top+\sum_{i=1}^{q-2}X^i p_i^\top\) must vanish on \(K_t\) at
\(\lambda_t\). The dual of this finite linear system says that
\(e_u^\top\sum_t z_t=0\) for every \(z_t\in K_t\) satisfying
\(\sum_t\lambda_t^i z_t=0\) for \(1\le i\le q-2\). Let
\(L(X)=\prod_t(X-\lambda_t)\) and interpolate the polynomial \(h\) with
\(\deg h<q\) and \(h(\lambda_t)=\lambda_tL'(\lambda_t)z_t\). For
\(\deg f<q\), the leading-coefficient formula from Lagrange
interpolation is \[
 [X^{q-1}] f=\sum_t f(\lambda_t)/L'(\lambda_t).
\] Applied to \(h\) this kills its top coefficient by the first moment.
As successive coefficients vanish, apply the formula to
\(Xh,X^2h,\ldots\). The \(q-2\) moments are equivalent to
\(\deg h\le1\). Conversely the same formula gives each moment from the
degree bound. The polynomial \(Ch\) has degree less than \(q\) and
vanishes at all \(q\) nodes, hence is zero. The interpolation formula at
zero gives \(h(0)=-L(0)\sum_t z_t\). The transformation is invertible on
the specified families and \(L(0)\ne0\). For \(q=2\) the moment list is
empty and \(\deg h<2\) already gives the claim. The output rows may be
solved independently. \(\square\)

Fix an output \(u=(a_u,b_u)\) in Lemma 4. At any position \(t\) with
\(a_u(t)=b_u(t)\), we have \(u\in G_{u,t}\), so the test imposes
\(h_u(\lambda_t)=0\). Two such positions give distinct roots of the
affine polynomial \(h_u\), forcing \(h_u\equiv0\); for that output the
conclusion of Lemma 4 is automatic. Call an index \emph{hard} if its
two words match in at most one position. For the lower bound we may
restrict the tests to vectors supported on hard column indices of
an arbitrary \(C\). In the upper construction every other column of
\(C\) is zero.

Write \(\bar a\) for the bitwise complement of \(a\) and \(a\oplus e_t\)
for flipping slot \(t\). Let \(v_a\) be the physical column of \(C\)
indexed by \((a,\bar a)\) and let \(c_{a,s}\) be its column indexed by
\((a,\bar a\oplus e_s)\). Put \(H=\operatorname{im} C\), \(m=\dim H\)
and \[
 N_a=\langle c_{a,1},\ldots,c_{a,q}\rangle,\qquad
 W_a=N_a+\langle v_a\rangle,\qquad \ell_t=1-X/\lambda_t.
\] Distinct \(\ell_t\) are linearly independent in pairs.

All outputs with second word \(\bar a\) have the same known sets. Their
coordinate constraints in Lemma 4, before imposing \(Ch(X)=0\), allow
affine test coordinates arbitrary in \(E\) on first word \(a\),
multiples of \(\ell_t\) on first word \(y=a\oplus e_t\), and zero on any
other first word. Thus the coefficient matrix of \(Ch(X)\) has columns
\(z\otimes1\) and \(z\otimes X\) for every hard physical column \(z\) at
\(a\), together with \(z\otimes\ell_t\) for every hard physical column
\(z\) at \(a\oplus e_t\). Since \(\ell_t(0)=1\), the coefficients tested
by Lemma 4 belong to the \(q+1\) columns \[
           v_a\otimes1,\qquad
           c_{a\oplus e_t,t}\otimes\ell_t\quad(1\le t\le q).
 \tag{3}
\] For feasibility, these columns must be independent modulo the span of
all other columns in the coefficient matrix. Indeed a relation with any
nonzero designated coefficient violates the corresponding output test.
For the construction supported on hard columns, this property is also
sufficient, since zero columns cannot cancel nonzero columns.

For each \(a\), the defining \(q+1\) vectors of \(W_a\) are independent.
A relation supported on first word \(a\) can be made a constant test for
center \(a\), forcing its \(v_a\) coefficient to be zero. For each
\(s\), multiply the same relation by \(\ell_s\) and apply the test for
center \(a\oplus e_s\) to force its \(c_{a,s}\) coefficient to be zero.
Hence \(\dim N_a=q\) and \(\dim W_a=q+1\).

\subsubsection{4.1 The lower bound and exceptional cube
edges}\label{the-lower-bound-and-exceptional-cube-edges}

The \(q\) neighbor columns in (3) remain independent modulo
\(W_a\otimes E\), because this latter subspace is spanned by the center
columns including the independent designated center column. The quotient
has dimension \(2(m-q-1)\), giving \(m\ge q+1+k\) for
\(k=\lceil q/2\rceil\).

Assume \(m=q+1+k\). For each center \(a\), divide \(H\otimes E\) by \[
             J_a=N_a\otimes E+\langle v_a\otimes X\rangle.
\] This removes only nondesignated center columns and has dimension
\(2q+1\). The designated columns remain independent modulo the
nondesignated neighbor columns. The images of all the latter therefore
span a space of dimension at most \[
             2m-(2q+1)-(q+1)=2k-q.
 \tag{4}
\] The value in (4) is zero for even \(q\) and one for odd \(q\).

At neighbor \(y=a\oplus e_t\) consider a nondesignated vector
\(z\in\{v_y\}\cup\{c_{y,s}:s\ne t\}\). If \(z\in W_a\mathbin{\backslash}N_a\),
then the class of \(z\otimes\ell_t\) modulo \(J_a\) is a nonzero
multiple of the designated center class \(v_a\otimes1\), an
impossibility. If \(z\notin W_a\), its image in the further quotient
\((H/W_a)\otimes E\) is a nonzero member of
\((H/W_a)\otimes\langle\ell_t\rangle\). For distinct \(t,t'\) the latter
subspaces intersect trivially. In view of (4), at most \(2k-q\)
directions at \(a\) can have a nondesignated neighbor vector outside
\(N_a\).

Call an incidence of a vertex \(a\) with an edge in direction \(t\)
exceptional when such a vector exists. At an incidence that is not
exceptional the \(q\) independent vectors \(v_y,c_{y,s}\) for \(s\ne t\)
all lie in \(N_a\), and they span it. If neither endpoint of an edge
were exceptional, we would have \(v_y\in N_a\subseteq W_y\) and
\(v_a\in N_y\subseteq W_a\). This would give \(W_a=N_a+N_y=W_y\). The
designated neighbor column would then belong to \(W_a\otimes E\),
contradicting its independence modulo that subspace. Each cube edge must
therefore have at least one exceptional incidence. The cube has
\(q2^{q-1}\) edges. For even \(q\) there are no exceptional incidences.
For odd \(q\ge3\) there are at most \(2^q<q2^{q-1}\). This contradiction
excludes \(m=q+1+k\) and proves \(m\ge q+k+2\).

\subsubsection{4.2 A binary-valued optimum
encoder}\label{a-binary-valued-optimum-encoder}

Use measurement coordinates \(e_1,\ldots,e_q\), \(f_1,\ldots,f_k\),
\(g_0,g_1\) and define the columns by \[
        c_{a,t}=e_t+a_t f_{\lceil t/2\rceil},
        \qquad
        v_a=g_{|a|\bmod2}.
 \tag{5}
\] All remaining physical columns are zero. The vectors at \(a=0\)
include the \(e_t\). Flipping the appropriate bit produces every \(f_i\)
by subtraction, and both word parities occur. This gives rank \(q+k+2\).

Fix center \(a\). In the span of the \(e\)- and \(f\)-coordinates the
\(f_i\) are independent modulo \(N_a\), by inspection of the \(e_t\)
coordinates. Project the affine coefficient matrix onto that quotient
tensored with \(E\). All center columns vanish. At neighbor
\(y=a\oplus e_t\) the vectors \(c_{y,s}\) for \(s\ne t\) agree with
\(c_{a,s}\). The surviving column has image \[
        (1-2a_t)\,\overline f_{\lceil t/2\rceil}\otimes\ell_t.
\] Each \(f_i\) occurs at most twice and the two associated polynomials
are independent. The scalar \(1-2a_t\) is either \(1\) or \(-1\), and
stays nonzero in characteristic two. Thus all \(q\) designated neighbor
coefficients vanish in any relation. On the separate \(g\)-coordinates,
every neighbor has parity opposite to \(a\). The coefficient of
\(g_{|a|\bmod2}\otimes1\) forces the designated center coefficient to
vanish. This proves (3) and completes the proof of Theorem 1.

For rational labels, the measurements in (5) also determine \(P\)
algorithmically. For a given output row, fix its constant term and
introduce unknowns for its remaining polynomial coefficients and the
combinations of the rows of \(C\) at each node. At every unavailable
input coordinate impose equality between the polynomial value and the
measurement combination. Lemma 4 and the preceding verification
guarantee that these rational linear systems are consistent. Gaussian
elimination constructs \(P\) in deterministic time polynomial in \(D\),
\(q\) and the bit length of the labels. After denominators are cleared,
the sizes of the minors bound the bit length needed by fraction-free
elimination.
The Lean sources also give a direct construction for this encoder. It
generates the coefficient computations from the slot count, evaluates
them on binary encodings of the rational labels, and writes every
coefficient of the resulting matrix within a polynomial bit budget.

\subsection{5. The exact cost of keeping the swap orbits
separate}\label{the-exact-cost-of-keeping-the-swap-orbits-separate}

The coordinate swaps exchange \(a_t\) and \(b_t\) independently. Their
orbits are specified by the unordered pair of letters at every slot. An
orbit with \(h\) unequal pairs is a Boolean cube of dimension \(h\). At
an active slot, the locally accessible points have the opposite
orientation to the output. At an inactive slot the entire orbit is
known. Call \(C\) an orbitwise direct sum if its row space is a direct
sum of subspaces each supported on one orbit. No constraint is placed on
those subspaces or on the coupling in \(P\).

\textbf{Proposition 5.} Over \(\mathbb Q\), for every \(r\ge2\) the
least orbitwise rank at degree zero is \((r^{2q}-r^q)/2\). For \(r=2\)
and \(q\ge2\), the least orbitwise rank at degree \(q-2\) is \(2q+2\).

\emph{Proof at degree zero.} In any nontrivial orbit, fix one active
slot and one orientation. Every input supported on this orientation has
zero side information for outputs on that orientation. Those outputs
must recover all \(2^{h-1}\) coordinates, forcing rank at least
\(2^{h-1}\) on the orbit. Pairing a point with its antipode and
transmitting each pair sum attains that rank. At any active slot the
partner can be read, and the point itself is available at each inactive
slot. The \(r^q\) singleton orbits contribute zero. Adding ranks proves
the first formula.

\emph{Lower bound at degree \(q-2\).} For \(r=2\), there are \(2q\)
orbits with exactly one inactive position and one all-active orbit. Each
of the former requires at least one measurement, since the polynomial on
a diagonal entry with no local measurement would otherwise have \(q-1\)
roots and intercept one. On the all-active orbit a zero measurement
column would similarly force \(q\) roots. If its rank were one and every
column nonzero, two adjacent physical points would give a relation
between their columns with both coefficients nonzero. Take an input
supported on these two points as prescribed by the relation. Its side
information is zero for one of the two outputs at their \(q-1\) common
orientations. The output row applied to the input would have those
\(q-1\) roots and nonzero intercept. This is impossible at degree
\(q-2\). Since an orbitwise direct sum charges the distinct orbits
separately, at least \(2q+2\) rows are necessary.

\emph{Attainment at degree \(q-2\).} On an orbit of dimension
\(h\le q-2\), take the identity on the orbit times
\(\prod_{t\text{ active}}(1-X/\lambda_t)\). For \(h=q-1\) measure the
constant Walsh row. For \(h=q\) measure the constant and full-parity
Walsh rows. To verify the resulting degree directly, index the
characters by subsets \(S\) of the active slots. The retained weights
form the interval \([L,U]=[1,h]\) or \([1,h-1]\), respectively.
Interpolate \(R(j)=(-1)^j\) on that interval by a polynomial of degree
\(U-L\). Give a retained character the multiplier \[
        p_S(X)=(-1)^{|S|}
          R\left(\sum_{i\in S}(1-X/\lambda_i)\right),
\] and use multiplier one on the measured characters. This
Walsh-diagonal matrix has constant term \(I\) and degree at most
\(q-2\). At an active slot \(t\), the multipliers of two retained
characters indexed by \(S\) and \(S\mathbin{\triangle}\{t\}\) have
opposite values at \(X=\lambda_t\). For a kernel input supported on the
unknown orientation, the two Fourier coefficients differ by the same
sign that relates the two characters at the output. Their contributions
cancel. If one character in a pair is measured, the two coefficients
both vanish. This verifies recovery on the unknown orientation, and
hence from the known orientation and the measurements. Inactive slots
already reveal the full orbit. This interpolation is the
retained-interval method from {[}1, 4{]}. The stated ranks follow.
\(\square\)

For \(r=2\) the orbitwise-to-unrestricted ratio at degree zero is
\(3\cdot2^q/[2(2^q+1)]\), tending to \(3/2\). At degree \(q-2\) it is
\((2q+2)/(q+\lceil q/2\rceil+2)\), tending to \(4/3\). These strict
inequalities for \(q\ge2\) show precisely where an optimal encoder must
couple different unordered-pair orbits.

\subsection{6. Verification, scope and open
cases}\label{verification-scope-and-open-cases}

Exact rational elimination checked the total-color construction at every
receiver on the full binary arrays for \(q=1,2,3\), with ranks
\(1,5,21\). It also checked all one-slot interval decoders for
\(2\le r\le8\). For (5), two calculations used the full arrays at
\(q=2,3,4\), with labels \(2,5,8,11\) truncated to the required number
of nodes. One eliminated the entrywise polynomial unknowns in the
original row equations. The other formed the divisibility matrix in
Lemma 4. Both certified all output rows at ranks \(5,7,8\). Removing a
row made half of the second-word blocks infeasible in each case.
Lowering the degree also made the tested constructions at \(q=3,4\)
fail. The accompanying Lean development formalizes the binary endpoints,
the general alphabet construction and its even-alphabet limit, the
receiver-order result, and the comparison with orbitwise encoders.
The repository includes a theorem index and verification logs.

The sender in these results observes the complete label array.
Determining the cost under distributed-sender restrictions requires
specifying which transmitted rows each sender can compute. The rank here
also concerns the full auxiliary model before physical variable
identification. Theorems 1 and 2 make no assertion about a SAT running
time. Two open cases are the limiting degree-zero cost for a fixed odd
alphabet larger than two and the optimum at general positive degree.

\textbf{Contribution and disclosure.} The paired polynomial recovery
model, the order-based zero-degree converse, standard matching
decompositions, and the moment and interval interpolation tools have the
antecedents identified above. The contributions developed here are the
proof of optimality of the proposed receiver order, the central
total-color encoders with their finite and limiting consequences, and
the exact degree-\(q-2\) rank with the cube-edge lower bound and the
binary-valued encoder. Chuhang Ding set the research questions and direction and proposed conjectures. AI assistants and agents contributed to proofs, tests, Lean formalization, review, and writing. Human peer review and a
determination of public priority remain open.

\subsection{References}\label{references}

{[}1{]} OpenAI. \emph{The Exponential Time Hypothesis is False}.
Revision of 25 October 2026, §4.1, Lemmas 4.1--4.3.
\href{https://raw.githubusercontent.com/openai/math/main/preprints/The-Exponential-Time-Hypothesis-is-False-October-25-2026/build/main_full_version.tex}{Primary
text}.

{[}2{]} Nicholas Woolsey, Rong-Rong Chen and Mingyue Ji. \emph{A
Combinatorial Design for Cascaded Coded Distributed Computing on General
Networks}. IEEE Transactions on Communications \textbf{69} (2021),
5686--5700. Theorems 2--3, Remark 5 and Appendices C--D.
\href{https://arxiv.org/pdf/2008.00581}{Author manuscript}.

{[}3{]} Prasad Krishnan, Lakshmi Natarajan and V. Lalitha. \emph{An
Umbrella Converse for Data Exchange: Applied to Caching, Computing, and
Shuffling}. Entropy \textbf{23} (2021), 985, §5.
\href{https://doi.org/10.3390/e23080985}{Journal article}.

{[}4{]} Vaibhav Tulsyan. \emph{What the ETH refutations do not give:
barriers for the counting engines, with machine-checked lemmas}.
November 2026, with companion note on the parity barrier, §§5.1--5.2.
\href{https://hexagonmath.org/pdf/2611.00149v1}{Paper} and
\href{https://github.com/xennygrimmato/eth-hamiltonicity/blob/946c49993b78aa163fcac3f8cb78ce351dd5dcc6/research/permanent/parity-barrier.md}{versioned
companion}.

\end{document}
